\begin{abstract}
Holography and the Riemann hypothesis are missing the same kind of object. The pair correlation of the zeros of
$\zeta$ follows GUE statistics; so does the boundary side of low-dimensional holography,
where the bulk of JT gravity is recovered only as an average over an ensemble of theories.
An ensemble is what one writes for want of an individual---the Hermitian operator of
Hilbert--P\'olya in one case, a single microstate in the other---and an average is a fact
about the description, not about what exists. Based on this reading, and on other known
connections between the primes and black holes, and between the odd zeta values and
string and particle amplitudes, we propose that the obstruction in de~Sitter is
\emph{epistemological}---a circumstance of mathematical circularity---rather than
\emph{ontological}---a matter of vacuum selection.

At this intersection the obstruction takes one form: how to describe the object without
referring to the object, which the $\F_1$ program seeks to circumvent in arithmetic and
the Hilbert--P\'olya problem poses for the operator. We trace it to the non-paradoxical
self-reference of binary language, in the precise sense of the Lawvere fixed-point theorem
and the Yanofsky diagonal $g(t)=\alpha(f(t,t))$, as one of the most studied cases of
self-referential circularity; and we analyse the strategies for confronting this
circularity and their relation to the $\F_1$ program for proving the Riemann
hypothesis---particularly Yuri Manin's proposed use of the same noncommutative tori that
M-theory uses, but to approach that program.

From our own use of this torus we construct a string cocone, from which we later derive
previously proposed criteria for holography in de~Sitter, as well as the operator we
propose as the microstate.
\end{abstract}
